\section{Literature Synthesis and Open Challenges}
\label{sec:literature-synthesis}
\label{sec:literature-synthesis-open-challenges}

The literature on stenotic hemodynamics is best read as a set of coupled
trade-offs rather than as a ladder from simple to complex models. Increasing
dimensional fidelity can expose local velocity gradients, wall interaction, and
geometry-specific flow features, but it also increases the number of boundary,
material, geometric, and numerical assumptions that must be identified. Reduced
models remove some of that burden and make system-level studies feasible, but
they move physical effects into closure laws, coupling conditions, and
observation maps. The open challenge is therefore not to choose one universal
model tier. It is to decide which tier can answer a specified question with
identifiable parameters, reproducible numerics, and a declared observable.

\subsection{Fidelity, Identifiability, and Observables}
\label{subsec:synthesis-fidelity-identifiability-observables}

Three-dimensional CFD and FSI formulations retain the most direct representation
of local flow and wall kinematics, while 1D and 0D models deliberately compress
that information into axial or lumped variables
\cite{FormaggiaEtAl2007FSICoupling,QuarteroniVeneziani2003GeometricalMultiscale,ShiEtAl2011BloodFlowModelReview}.
This compression is not only a numerical convenience. It is also an
identifiability choice. A highly resolved forward model can depend on material
properties, prestress, inlet histories, outflow impedances, geometry state, and
wall motion that are rarely all observed with matching certainty. Conversely, a
1D model may have fewer state variables but can still be underdetermined if wall,
profile, friction, rheology, and boundary closures are treated as adjustable
without independent evidence. Parameter-identifiability work for hemodynamics
PDE models makes this tension explicit: adding modeling detail is valuable only
when the data and outputs can distinguish the resulting parameters
\cite{Colebank2025HemodynamicsIdentifiability}.

This is why observation operators are part of the model, not post-processing
decorations. A 1D area--flow solution does not natively contain the same
observable as a resolved 3D velocity field, and a pressure-drop, flow-rate,
section-mean velocity, wall-shear, or FFR-style quantity each imposes a different
comparison contract. Multiscale coupling work already treats interface
quantities and compatibility conditions as mathematical structure
\cite{FormaggiaEtAl2007FSICoupling,QuarteroniVeneziani2003GeometricalMultiscale};
diagnostic 1D--3D comparison needs the same discipline at the output side.

\subsection{Uncertainty in Boundaries, Walls, Rheology, and Geometry}
\label{subsec:synthesis-closure-geometry-uncertainty}

Boundary-condition uncertainty is one of the most persistent barriers to
interpretable hemodynamics. Inlet waveforms, outlet impedances, terminal
resistances, and coupled-domain histories can change the local solution as much
as the internal discretization. In a multiscale setting, these quantities are not
auxiliary inputs; they are the mechanism by which upstream and downstream
vascular beds communicate with the modeled stenosis
\cite{QuarteroniVeneziani2003GeometricalMultiscale,FormaggiaEtAl2007FSICoupling}.
A reduced stenosis result that does not declare its boundary approximation
therefore cannot be cleanly compared with a resolved result, even if both
models share a nominal geometry.

Wall and rheology uncertainty create a second closure layer. Wall stiffness and
thickness determine wave speed and area response; rheology and velocity-profile
assumptions determine dissipative behavior and momentum transport. In 1D models,
these effects enter through pressure--area laws, friction terms, effective
viscosity choices, and momentum-flux correction factors
\cite{ShiEtAl2011BloodFlowModelReview}. In FSI models, they enter through wall
material laws, structural boundary conditions, and the coupling algorithm. Both
settings face the same interpretive problem: when the wall or rheology is not
independently constrained, a discrepancy can be a physics signal, a closure
signal, or a compensation between uncertain choices.

Geometry is similarly double-edged. Patient-specific geometry is attractive
because it carries clinically meaningful morphology, but segmentation,
registration, current-versus-reference configuration, and wall-motion state can
dominate the modeling record. Idealized geometry removes many of those
ambiguities and exposes controlled severity and refinement parameters, but it
cannot by itself answer patient-specific questions. The useful distinction is
not realistic versus artificial. It is whether the geometry serves the claim:
patient-specific geometry for patient-specific interpretation, and idealized
geometry for isolating numerical, closure, and comparison behavior.

\subsection{Consistency Across Scales and Discretizations}
\label{subsec:synthesis-consistency-verification-validation}

Multiscale consistency requires more than matching units at an interface. The
model tier, the boundary representation, the wall law, and the output operator
must all describe compatible quantities. This is the lesson connecting
geometrical multiscale models with 1D benchmark practice: numerical results are
interpretable only when the mathematical problem, coupling assumptions, and
solver realization are stated together
\cite{QuarteroniVeneziani2003GeometricalMultiscale,BoileauEtAl2015Benchmark}.
Open 1D solvers such as openBF further emphasize that reproducible workflows,
transparent inputs, and inspectable finite-volume records are part of the
scientific object, not merely software packaging
\cite{BenemeritoEtAl2024OpenBF}.

Equilibrium preservation is a particularly sharp consistency test. A continuous
stenotic balance law may admit a rest or steady state through cancellation
between flux gradients, geometry terms, wall response, friction, and boundary
data. A discrete method that does not preserve that state can generate motion
from the numerical representation itself. Well-balanced methods for 1D blood
flow make this point directly and raise the standard for stenosis computations:
manufactured solutions and grid studies are useful, but an equilibrium problem
needs an equilibrium-preservation audit
\cite{PimentelGarciaEtAl2025HighOrderWellBalanced}. The broader verification
and validation gap is that many studies report plausible hemodynamic outputs
without enough evidence to separate implementation error, discretization error,
closure uncertainty, cross-model mismatch, and external validation against an
accepted target.

\subsection{Reproducibility and Emerging Learned Methods}
\label{subsec:synthesis-reproducibility-learned-methods}

Reproducibility is the practical bridge between modeling ambition and reviewable
evidence. Benchmark studies help by fixing common cases and comparison
quantities, but stenosis studies also need persisted geometry definitions,
boundary histories, closure parameters, solver settings, diagnostics, and output
maps \cite{BoileauEtAl2015Benchmark,BenemeritoEtAl2024OpenBF}. Without those
records, a disagreement between two models is hard to assign to physics,
parameters, numerics, or data handling.

Learned methods do not remove this requirement. Physics-informed neural networks
offer a framework for solving and inverting PDE-constrained problems
\cite{RaissiEtAl2019PhysicsInformedNeuralNetworks}, and hemodynamics reviews
now treat PINNs as an active direction for vascular modeling
\cite{YuEtAl2026PINNHemodynamicsReview}. At the same time, gradient-pathology
analyses show that training dynamics can be a limiting numerical issue rather
than a transparent solver detail
\cite{WangEtAl2021PINNGradientPathologies}. Operator-learning methods such as
DeepONet and Fourier neural operators shift the ambition from one solution to a
map between function spaces \cite{LuEtAl2021DeepONet,LiEtAl2021FNO}, and recent
stenosed-vessel work shows why such models are appealing for fast parameterized
flow prediction \cite{VelikorodnyEtAl2025DeepOperatorStenosed}. Their open
challenge is the same synthesis problem in a new form: the training distribution,
geometry family, boundary inputs, wall and rheology assumptions, loss
normalization, and validation observables must be declared before speed or fit
can be interpreted as hemodynamic evidence.

\subsection{Role of the Idealized Stenosis Case Study}
\label{subsec:synthesis-idealized-stenosis-case-study}

An idealized stenosis case study is useful within this review because it turns
the literature trade-offs into a worked audit. It fixes a geometry family,
closure contract, boundary approximation, discretization, and observation
operator tightly enough that the consequences can be inspected. The same setup
can ask whether a reduced model preserves a rest equilibrium, whether a velocity
comparison is being made through a declared cross-sectional operator, whether
the recorded artifacts are sufficient to reproduce the numerical claim, and
which uncertainties remain outside the case.

The accompanying \texttt{StenosisHemodynamics} implementation therefore has a
methodological role. The case removes segmentation and clinical-matching
ambiguity so that closure dependence, boundary sensitivity, equilibrium
preservation, verification evidence, and cross-dimensional observability can be
examined in a controlled setting. That is exactly the role a worked example
should play in a review-led manuscript: it demonstrates how the field's open
challenges appear in one auditable stenosis problem, while patient-specific
interpretation remains the domain of studies with matched patient geometry,
boundary data, wall state, and validation targets.
